\chapter{Metadata providers}
\label{ch:providers}

When you scan an ISBN or an EAN-13, mybibli does not look the title
up in a local catalog — it queries one or more external
\emph{metadata providers}\index{metadata providers} until one
returns a clean answer. This chapter explains the chain, where to
get API keys, and how to rotate them safely.

\section{The provider chain}

mybibli walks the providers in the order best-suited to the guessed
media type:

\begin{longtable}{@{} l l p{4.5cm} c @{}}
\toprule
\textbf{Provider} & \textbf{Media types} & \textbf{Notes} & \textbf{Key?} \\
\midrule
\endhead
BnF\index{BnF}                  & Books         & French
Bibliothèque nationale catalogue; best for French-language titles. & no \\
Open Library\index{Open Library} & Books         & Volunteer
catalogue; broad English / multilingual coverage. & no \\
Google Books\index{Google Books} & Books         & Commercial; good
for trade non-fiction. & optional \\
MusicBrainz\index{MusicBrainz}  & Audio         & Volunteer music
metadata; very deep on CDs / LPs. & no \\
OMDb\index{OMDb}                & Films, series & Free tier; rate-limited. & required \\
TMDb\index{TMDb}                & Films, series & Free with attribution. & required \\
BDGest\index{BDGest}            & Comics, BD    & FR-language comic
catalogue (web-scraped). & no \\
\bottomrule
\end{longtable}

The first provider to return a \emph{clean} record wins. ``Clean''
means: title present, at least one contributor, and a year. Partial
matches are recorded but mybibli keeps walking the chain until it
finds a clean one or runs out of providers.

\section{API keys}

Three providers require an API key:

\subsection*{Google Books}

Optional but raises the daily request budget significantly.

\begin{enumerate}
  \item Sign in to \url{https://console.cloud.google.com/}.
  \item Create a new project (or reuse one).
  \item Enable the \textbf{Books API}.
  \item Create an \textbf{API key} under \emph{APIs \& Services →
        Credentials}.
  \item Paste it into \texttt{/admin > System > Metadata Providers}
        in mybibli.
\end{enumerate}

\subsection*{OMDb}

\begin{enumerate}
  \item Visit \url{https://www.omdbapi.com/apikey.aspx}.
  \item Pick the free tier (1000 requests/day).
  \item Confirm via the activation email you receive.
  \item Paste the key into mybibli.
\end{enumerate}

\subsection*{TMDb}

\begin{enumerate}
  \item Sign up at \url{https://www.themoviedb.org/}.
  \item Open \emph{Settings → API} and request a key for
        \emph{Personal / non-commercial use}.
  \item Paste the key into mybibli.
\end{enumerate}

\section{Where the keys live}

Keys travel through two layers:

\begin{enumerate}
  \item \textbf{Environment variables}\index{env vars!provider keys}
        (\texttt{GOOGLE\_BOOKS\_API\_KEY}, \texttt{OMDB\_API\_KEY},
        \texttt{TMDB\_API\_KEY}). These are migrated \emph{once} on
        boot into the \texttt{settings} table — they are essentially
        a deployment-time bootstrap.
  \item \textbf{The \texttt{settings} table}, exposed through
        \texttt{/admin > System > Metadata Providers}. This is the
        live source of truth at request time. Rotating a key from
        the UI takes effect on the very next metadata fetch.
\end{enumerate}

The implication is asymmetrical:

\begin{itemize}
  \item If you set a key only in the environment, the first boot
        copies it into the database and afterwards the database
        wins.
  \item If you set a key only in the UI, the environment is never
        consulted again.
  \item If you set both and then \emph{Clear} the key from the UI,
        the database value goes empty — but on the next boot, the
        environment value is migrated back in. To make the UI clear
        durable, also remove the environment variable from
        \texttt{.env} / \texttt{docker-compose.yml}.
\end{itemize}

\section{Provider health}

\texttt{/admin > Health} pings every registered provider every five
minutes and reports their status as a green check or a red cross.
A red cross does not always mean the provider is down — it can also
mean the API key is missing or invalid. The Health tab includes the
last-checked timestamp and the HTTP status code for diagnostics.

\section{Rate limits and back-pressure}

mybibli applies a small per-provider rate limit (a couple of
requests per second). Bulk-cataloging a hundred volumes in a row
will not trigger remote 429 responses. If a provider does return
429, the request is logged and the next provider in the chain is
consulted instead — the catalog flow never blocks waiting for a
rate-limited provider.

\begin{warning}
Never commit a real API key to git. The \texttt{.env} file is in
\texttt{.gitignore} for this reason; \texttt{.env.example} carries
only placeholders. If a key leaks, rotate it at the provider
console immediately — pasting a new key into mybibli takes effect
within seconds.
\end{warning}
